%% polyglossia: this is rtl.tex with _preamble-polyglossia, the body
%% kept identical so that a_rtl asserts on it unchanged.  A difference
%% between the two is a difference polyglossia (and, under XeLaTeX,
%% bidi) makes.  Keep the bodies in step.
%%
%% _preamble.tex with polyglossia loaded BEFORE linguexx, Hebrew among its
%% languages.  Under XeLaTeX that pulls in bidi, which redefines boxes,
%% lists and footnotes wholesale -- the realistic setting for a right-to-
%% left gloss, and one no other preamble here loads.  XeLaTeX and LuaLaTeX
%% only: polyglossia is fontspec's.
\documentclass[11pt]{article}
\usepackage[a4paper,margin=25mm]{geometry}
\usepackage{fontspec}
\usepackage{polyglossia}
\setdefaultlanguage{english}
\setotherlanguages{hebrew,german}
\newfontfamily\hebrewfont{FreeSerif.otf}[Script=Hebrew]
\usepackage{linguexx}
\pagestyle{empty}

%% \GlossRTL on the page: columns from the right, a short grid beside its
%% number, a long one wrapping from the right, the number and a judgment
%% mark on the left, and \GlossRTLOff back to normal.
%%
%% Measured on the Latin tiers only: pdftotext reads Hebrew in visual
%% order under XeTeX, and the Latin words are sentinels.  FreeSerif is
%% TeX Live's (gnu-freefont), by file name so that both engines find it.
\newfontfamily\hebfont{FreeSerif.otf}[Script=Hebrew]
\begin{document}
\GlossTierLang{1}{he}
\GlossTierFont{1}{\hebfont}
% tier 2 is a transliteration: its script subtag makes it left to right,
% so its words read as typed (a right-to-left tier 2 would reverse them)
\GlossTierLang{2}{he-Latn}
\GlossRTL

% (1) a short grid.  Columns run leftwards; in each column the cells sit
% against its right edge (RBONEXX is wider than RAONE on purpose); and
% the grid starts at the text margin, beside its number, where the
% translation does.
\ex. \glll הילד אכל את \\
           RAONE RATWO RATHREE \\
           RBONEXX RBTWO RBTHREE \\
\glt RTRANSONE starts at the text margin.

% (2) a grid too long for one line: LWAA is the first word, so the
% rightmost of the first line, and both lines end at the right margin.
\ex. \gll הילד אכל את התפוח הגדול שקנינו אתמול בשוק של העיר העתיקה ליד הנמל הישן \\
          LWAA LWAB LWAC LWAD LWAE LWAF LWAG LWAH LWAI LWAJ LWAK LWAL LWAM LWAN \\
\glt RTRANSTWO

% (3) sub-examples, and a judgment mark: on the left, where it hangs in
% any example (provisional -- see doc/DEFERRED-DECISIONS.md).
\ex. \a. \gll הילד אכל \\ SUBONE SUBTWO \\
     \b. *\gll אכל הילד \\ JUDONE JUDTWO \\

% (5) the grid is measured before it is placed, so it is set twice: a
% morpheme mismatch is still reported once, not once per setting.
\GlossMorphAlign
\ex. \glll הילד \\ MMA-mb \\
          mx-my-mz \\
\GlossMorphAlignOff

% (4) off again: left to right.
\GlossRTLOff
\GlossTierFont{1}{\textnormal}\GlossTierLang{1}{de}
\ex. \gll Der Hund \\ OFFONE OFFTWO \\

% (6) a script line, right to left whether or not the grid is: its first
% word (the one with the dalet) stands to the right of its second (the one
% with the kaf).  Under LuaTeX pdftotext gives single letters, under XeTeX
% whole words; the comparison holds for both.
\ex. \glscript{he}{\hebfont הילד אכל}
\gll SCRA SCRB \\ scra scrb \\

\end{document}
